\documentclass[11pt,a4paper]{article}
\usepackage[margin=27mm]{geometry}
\usepackage{amsmath,amssymb,amsthm,mathtools}
\usepackage{fontspec}
\setmainfont{TeX Gyre Pagella}
\setsansfont{TeX Gyre Heros}
\setmonofont[Scale=0.82]{Latin Modern Mono}
\usepackage{microtype}
\usepackage{xcolor,booktabs,tabularx,longtable,enumitem}
\definecolor{inkblue}{HTML}{244860}
\usepackage[unicode,colorlinks=true,linkcolor=inkblue,urlcolor=inkblue,citecolor=inkblue]{hyperref}
\usepackage{bookmark}
\usepackage{fancyhdr}
\pagestyle{fancy}\fancyhf{}
\fancyhead[L]{\small\sffamily Cluster skeleton geometry}
\fancyhead[R]{}
\fancyfoot[C]{\small\thepage}
\setlength{\headheight}{15pt}
\setlength{\parskip}{4pt}
\setlength{\parindent}{0pt}
\setlist{itemsep=3pt,topsep=5pt}
\newtheorem{theorem}{Theorem}[section]
\newtheorem{lemma}[theorem]{Lemma}
\newtheorem{proposition}[theorem]{Proposition}
\newtheorem{corollary}[theorem]{Corollary}
\theoremstyle{definition}
\newtheorem{definition}[theorem]{Definition}
\newtheorem{remark}[theorem]{Remark}
\newcommand{\C}{\mathbb C}
\newcommand{\R}{\mathbb R}
\newcommand{\Z}{\mathbb Z}
\newcommand{\dd}{\mathrm d}
\newcommand{\im}{\operatorname{Im}}
\newcommand{\re}{\operatorname{Re}}
\newcommand{\Skel}{\operatorname{Skel}}
\newcommand{\supp}{\operatorname{supp}}
\newcommand{\lean}[1]{\texttt{\small\detokenize{#1}}}
\newcommand{\leanref}[2]{\par\smallskip{\small\sffamily Lean correspondence: \lean{#1}; \lean{#2}.}\par\smallskip}
\emergencystretch=2em

\title{Liouville geometry and the full skeleton\\[5pt]\large A geometric proof for unimodular acyclic cluster models}
\author{Cluster skeleton formalization project}
\date{21 October 2026}
\begin{document}
\maketitle
% Body only. Requires amsmath, amssymb, amsthm, mathtools, booktabs,
% longtable, hyperref; theorem/lemma/proposition/definition/remark environments.
% Use an 11pt article and ordinary margins; no external bibliography is needed.

\begin{abstract}
We construct a smooth exact symplectic structure on the full exchange
variety of an integral unimodular skew-symmetric seed whose mutable quiver
is acyclic. Frozen coordinates are retained as invertible variables.
The construction extends the imaginary part of the logarithmic two-form
from the initial torus, has a Liouville vector field complete in both time
directions, and identifies its full trapped set with an explicit compact
coordinate locus. The proof treats all simultaneous-zero strata and uses
the same vector field and complete flow throughout. This note gives the
mathematical argument for the geometric statement formalized in the
accompanying Lean project.
\end{abstract}

\section{The exchange variety and the precise statement}

Let $E=\{1,\ldots,d\}$, with a partition $E=M\sqcup F$ into mutable
and frozen indices. Let $B=(B_{ij})_{i,j\in E}$ be an integral matrix
satisfying
\[
 B_{ij}=-B_{ji},\qquad \det B=1.
\]
An entry $B_{ij}>0$ denotes arrows from $i$ to $j$, and the directed
graph induced on $M$ is assumed acyclic. Multiple arrows are allowed.
The determinant hypothesis concerns the \emph{whole} matrix, including
the frozen rows and columns. The empty index set is allowed, with the
usual conventions for empty products and the determinant of the empty
matrix.

For $i\in M$, define the exchange monomials
\[
 P_i(a)=\prod_{j\in E}a_j^{\max(B_{ij},0)},\qquad
 Q_i(a)=\prod_{j\in E}a_j^{\max(-B_{ij},0)}.
\]
Our space is the actual coordinate locus
\begin{equation}\label{eq:exchange}
 Y=\left\{(a,b)\in\mathbb C^E\times\mathbb C^M:
 a_i b_i=P_i(a)+Q_i(a)\ (i\in M),\quad
 a_f\ne0\ (f\in F)\right\},
\end{equation}
with its inherited topology. In particular, the frozen coordinates are
variables, rather than prescribed coefficients. Partner coordinates
$b_i$ occur only at mutable indices. The point $a_j=1$, $b_i=2$ shows
that the space is nonempty.

The initial torus is
\[
 T=\{p\in Y:a_j(p)\ne0\text{ for all }j\in E\}.
\]
Eliminating $b_i$ identifies $T$ with $(\mathbb C^*)^d$. On this torus,
put
\begin{equation}\label{eq:torusform}
 \Omega_T=\frac12\sum_{i,j\in E}B_{ij}
               \frac{da_i}{a_i}\wedge\frac{da_j}{a_j}.
\end{equation}
The factor $1/2$ compensates exactly once for the two ordered pairs
$(i,j)$ and $(j,i)$.

Independently of any differential form or flow, define the coordinate
cap
\begin{equation}\label{eq:cap}
 C=\left\{p\in Y:
 |a_f|=1\ (f\in F),\quad |a_i|\le1,\quad
 (1-|a_i|^2)b_i=0\ (i\in M)\right\}.
\end{equation}
Thus a partner must vanish when its original coordinate is strictly
inside its unit disk; no additional partner condition is imposed by
\eqref{eq:cap} when $|a_i|=1$. The exchange equations always remain in
force.

For a complete flow $\varphi:\mathbb R\times Y\to Y$, we use the
following definition:
\begin{equation}\label{eq:trapped}
 \operatorname{Skel}(\varphi)=
 \left\{p\in Y:
 \overline{\{\varphi_t(p):t\ge0\}}\text{ is compact in }Y\right\}.
\end{equation}
This definition uses the entire positive orbit and the topology of $Y$.
It does not assume a compact-domain presentation of the manifold.

\begin{theorem}\label{thm:main}
The space $Y$ admits a smooth real manifold structure of dimension $2d$
compatible with its inherited topology, for which the coordinate
inclusion in $\mathbb C^E\times\mathbb C^M$ is a smooth immersion.
The initial torus is dense. There is a smooth real one-form $\lambda$
such that:
\begin{enumerate}
\item $\omega=d\lambda$ is symplectic and
      $\omega|_T=\operatorname{Im}\Omega_T$;
\item the vector field $Z$ determined by $\iota_Z\omega=\lambda$ has a
      smooth flow $\varphi$ defined for all real times;
\item for this same flow,
      $\operatorname{Skel}(\varphi)=C$.
\end{enumerate}
\end{theorem}

The statement proved here is about smooth real forms and complete
Liouville dynamics. A holomorphic extension theorem is not a premise of
the construction. The proof also does not need a Weinstein structure,
finite-type conclusion, cap product decomposition, or categorical
identification.

\section{Integral node charts on the full coordinate locus}

A subset $I\subseteq M$ is \emph{independent} if $B_{ij}=0$ for all
$i,j\in I$. For a point $p\in Y$, write
\[
 I(p)=\{i\in M:a_i(p)=0\}.
\]

\begin{lemma}\label{lem:zeros}
The set $I(p)$ is independent for every $p\in Y$.
\end{lemma}
\begin{proof}
Suppose that the induced directed graph on $I(p)$ contains an edge.
Choose a source in a nontrivial connected component of that graph; this
is possible because the mutable graph is finite and acyclic. Let $i$
be this source. At least one outgoing neighbor has zero coordinate, so
$P_i(p)=0$. There is no incoming neighbor in $I(p)$, and every frozen
coordinate is nonzero; hence every factor in $Q_i(p)$ is nonzero.
Thus $Q_i(p)\ne0$. The equation $a_i b_i=P_i+Q_i$ becomes
$0=Q_i(p)$, a contradiction. Skew-symmetry then excludes both possible
signs of an entry between two zero indices.
\end{proof}

For an independent $I$, let $L=E\setminus I$ and set
\[
 U_I=\{p\in Y:a_j(p)\ne0\text{ for all }j\in L\}.
\]
For $i\in I$, every variable appearing in $P_i$ or $Q_i$ belongs to
$L$. Consequently $Q_i$ is a unit on $U_I$, and the expressions
\begin{equation}\label{eq:normalized}
 q_i=P_i/Q_i=\prod_{j\in L}a_j^{B_{ij}},\qquad
 v_i=b_i/Q_i,\qquad u_i=a_i
\end{equation}
are well-defined there. Their exchange equation is
\begin{equation}\label{eq:node}
 u_i v_i=1+q_i,\qquad q_i\ne0.
\end{equation}

\begin{lemma}\label{lem:lattice}
The rows $(B_{ij})_{j\in L}$, indexed by $i\in I$, extend to an
integral basis of $\mathbb Z^L$.
\end{lemma}
\begin{proof}
Since $\det B=1$, all rows of $B$ form an integral basis of
$\mathbb Z^E$. The span of the selected rows is therefore a direct
summand, and in particular is saturated. These rows have zero
$I$-coordinates by independence. Their span lies in the direct summand
$\mathbb Z^L\subseteq\mathbb Z^E$ and remains saturated there: if a
nonzero integer multiple of $x\in\mathbb Z^L$ belongs to the span,
saturation in $\mathbb Z^E$ already puts $x$ in that span. The selected
rows are independent. Their quotient in $\mathbb Z^L$ is thus a
finitely generated torsion-free abelian group, hence free. Lifting a
basis of the quotient gives the required completion.
\end{proof}

Fix such a basis completion. It gives Laurent coordinates
\[
 z=(q_i\ (i\in I),w_1,\ldots,w_{d-2|I|})
\]
on $(\mathbb C^*)^L$. The inverse change of variables also consists of
Laurent monomials with integral exponents. Let $F$ denote the matrix of
the forward change and $D=F^{-1}$ its integral inverse, so that
\[
 a_j=\prod_k z_k^{D_{jk}}\qquad(j\in L).
\]
There are $d-2|I|$ complementary coordinates because the selected
independent rows number $|I|$ inside a lattice of rank $d-|I|$.

Eliminate $q_i$ by $q_i=u_i v_i-1$. For mutable $j\notin I$, eliminate
the partner using $b_j=(P_j+Q_j)/a_j$. This produces an explicit
coordinate isomorphism
\begin{equation}\label{eq:chart}
 U_I\simeq
 \left\{(u_i,v_i)_{i\in I}:u_i v_i\ne1\right\}
 \times(\mathbb C^*)^{d-2|I|}.
\end{equation}
The remaining partners are recovered by $b_i=Q_i v_i$. Every denominator
used here is nonzero on the stated domain. In particular, this is an
integral Laurent change of coordinates, with no finite covering.

The right side of \eqref{eq:chart} is an open subset of a complex vector
space of dimension $d$, viewed as a real vector space of dimension
$2d$. The forward and inverse maps are continuous in the inherited
topology and smooth on their domains. On overlaps they give smooth
transition maps. Lemma~\ref{lem:zeros} ensures that a point belongs to
$U_{I(p)}$, so these charts cover $Y$. They define the claimed real
smooth structure. The coordinate inclusion is smooth; it is an
immersion because the smooth chart map, defined using the ambient
coordinates on a neighborhood, is a local left inverse of the
coordinate inclusion and therefore makes its derivative injective.

In a node chart, $T$ is obtained by also requiring $u_i\ne0$ for each
$i\in I$. The complement is a finite union of coordinate hyperplanes
inside an open domain. The torus is therefore dense in every node
chart, and hence dense in $Y$. We will use this density to glue actual
smooth forms and to extend identities of continuous tensors.

\section{The global real symplectic form}

We retain the entire residual block in the node expression of
\eqref{eq:torusform}. This avoids imposing a nondegeneracy assumption
on an auxiliary form.

Let $B_{LL}$ be the outside block of $B$, and put
\[
 R=D^{\mathsf T}B_{LL}D.
\]
The indices of $R$ are the $q$-indices and the complementary $w$-indices.
On the node domain, define
\[
 \eta_i=\frac{d(u_i v_i-1)}{u_i v_i-1}\quad(i\in I),
 \qquad \eta_k=\frac{dw_k}{w_k}\quad(k\text{ complementary}).
\]
Consider the complex-valued expression
\begin{equation}\label{eq:nodeform}
 \Omega_I=\sum_{i\in I}\frac{du_i\wedge dv_i}{u_i v_i-1}
       +\frac12\sum_{k,l}R_{kl}\eta_k\wedge\eta_l.
\end{equation}
All its coefficients are smooth on the node domain. On the initial
torus, independence splits \eqref{eq:torusform} as
\[
 \Omega_T=\sum_{i\in I}d\log a_i\wedge d\log q_i
       +\frac12\sum_{j,k\in L}B_{jk}\,d\log a_j\wedge d\log a_k.
\]
The first summand has the exact cancellation
\begin{equation}\label{eq:nodecancellation}
 d\log u_i\wedge d\log(u_i v_i-1)
 =\frac{du_i\wedge dv_i}{u_i v_i-1}.
\end{equation}
The inverse Laurent differential expresses the second summand by $R$.
Thus \eqref{eq:nodeform} agrees with \eqref{eq:torusform} on $T\cap U_I$.

Pull back $\operatorname{Im}\Omega_I$ by the actual node-chart
differential. On overlaps the resulting real two-forms agree on the
dense torus. Their coefficients are continuous, so they agree on the
whole overlap. They glue to a smooth real two-form $\omega$ on $Y$.
The identity $d\omega=0$ follows in the same way: it holds on $T$ by
the logarithmic formula, and $d\omega$ is a smooth, hence continuous,
tensor. This establishes a global real extension without using a
separately assumed holomorphic extension.

\begin{lemma}\label{lem:nondegenerate}
The form $\omega$ is nondegenerate at every point of $Y$.
\end{lemma}
\begin{proof}
Choose the chart indexed by the actual zero set $I=I(p)$. At its
center $u_i=0$ for all $i\in I$, so $q_i=-1$ and
\[
 \eta_i=-v_i\,du_i.
\]
After the integral outside change of basis, the matrix of the original
logarithmic form, in the order $(a_I,q_I,w)$, is
\begin{equation}\label{eq:block}
 \begin{pmatrix}
 0&1&0\\
 -1&R_{qq}&R_{qw}\\
 0&R_{wq}&R_{ww}
 \end{pmatrix}.
\end{equation}
It is invertible because it is congruent to the full invertible matrix
$B$. In particular, $R_{ww}$ is invertible. Indeed, if $R_{ww}x=0$,
then $(R_{qw}x,0,x)$ lies in the kernel of \eqref{eq:block}, and hence
$x=0$. This proof also covers a zero-dimensional complementary block.

Now let a complex tangent vector at the center be in the kernel of
$\Omega_I$. Test it first against vectors supported in the $v_i$
directions. The first term of \eqref{eq:nodeform} is
$-\sum_i du_i\wedge dv_i$ there, while the residual term has no $dv$
component. Thus the tested vector has zero $u$-components. Testing next
against the free $w$-directions and using the invertibility of $R_{ww}$
forces its $w$-components to vanish. Finally testing against the
$u_i$-directions forces its $v$-components to vanish. Therefore the
complex bilinear expression is nondegenerate at the center.

For a complex bilinear alternating form, nondegeneracy implies
nondegeneracy of its imaginary part as a real form: if a real tangent
vector $x$ satisfies $\operatorname{Im}\Omega_I(x,y)=0$ for every $y$,
testing also against $\mathrm i y$ gives
$\operatorname{Re}\Omega_I(x,y)=0$. Hence $x=0$. The chart differential
is a real tangent isomorphism, so the same conclusion holds for
$\omega_p$. Every point is the center of the chart indexed by its own
zero set; no density argument for nondegeneracy is required.
\end{proof}

\section{An explicit primitive on the torus}

On $T$, write
\[
 A_j=\log|a_j|,\quad \theta_j=\arg(a_j),\quad
 S_i=\sum_j B_{ij}A_j=\log|q_i|,\quad T_i=\arg(-q_i).
\]
Angles are circle-valued. Their differentials are globally defined on
the torus, and a periodic scalar function of $T_i$ is independent of
its representative. Set
\begin{equation}\label{eq:beta}
 \beta=\sum_{i,j} A_i B_{ij}\,d\theta_j.
\end{equation}
Expanding $d\log a_i=dA_i+\mathrm i\,d\theta_i$ and using
skew-symmetry gives $d\beta=\operatorname{Im}\Omega_T$.

We give concrete cutoff choices. Write
\[
 e(x)=\begin{cases}e^{-1/x},&x>0,\\0,&x\le0,\end{cases}
 \qquad \sigma(x)=\frac{e(x)}{e(x)+e(1-x)}.
\]
This is the smooth monotone transition used in Mathlib: $\sigma(x)=0$ for
$x\le0$, $\sigma(x)=1$ for $x\ge1$, and
$\sigma(1-x)=1-\sigma(x)$. Fix
\begin{equation}\label{eq:parameters}
 \varepsilon>0,\qquad \kappa\ge 2\varepsilon+1,\qquad\delta>0.
\end{equation}
The choice $(\varepsilon,\kappa,\delta)=(1,3,1)$ will suffice for the
existence theorem.

Define
\begin{align}
 \chi(x)&=\sigma\left(\frac{-2x-\varepsilon}{\varepsilon}\right),
 &\rho(s)&=\sigma(2-s^2),\label{eq:cutoffs}\\
 \chi_i(A)&=\prod_{\substack{j\in M\\B_{ij}\ne0}}
                     \chi(A_i-A_j).\label{eq:dominance}
\end{align}
Then $\chi=1$ on $(-\infty,-\varepsilon]$ and $\chi=0$ on
$[-\varepsilon/2,\infty)$. The function $\rho$ equals one for
$|s|\le1$ and has compact support. Only mutable neighbors enter
\eqref{eq:dominance}. In particular, frozen coefficients are not
fixed or compared by this cutoff.

For adjacent mutable indices $i,j$, $\chi_i$ and $\chi_j$ cannot
both be nonzero. Indeed nonvanishing would require both
$A_i-A_j<-\varepsilon/2$ and
$A_j-A_i<-\varepsilon/2$. This exclusion is the reason for using the
dominance factors: a bounded ratio $q_i$ alone does not detect whether
several of its factors are approaching zero.

Let
\[
 L=2(\kappa-2\varepsilon)>0,\qquad
 G(x)=\int_0^x\sigma(t)\,dt,
 \qquad h(A)=-L\,G\left(\frac{-A-2\varepsilon}{L}\right).
\]
The reflection identity for $\sigma$ gives $G(1)=1/2$, and direct
differentiation yields
\begin{gather}
 h\le0,\qquad 0\le h'\le1,\qquad |h(A)|\le|A|,
 \label{eq:hproperties}\\
 h(A)=0\quad(A\ge-2\varepsilon),\qquad
 h(A)=A+\kappa\quad(A\le-L-2\varepsilon),\label{eq:htail}\\
 A<h(A)\quad(A<0).
 \label{eq:hstrict}
\end{gather}
For the last inequality, if $A<-2\varepsilon$, use $0\le G(x)\le x$
for $x\ge0$ to obtain $h(A)\ge A+2\varepsilon>A$; in the other
negative range it is immediate. All positive-order derivatives of
$h$ are bounded. For order at least two this also follows from compact
support of the second derivative and its successive derivatives.

For the angular cutoff put $r=\min(\delta/2,\pi/4)$ and
\begin{equation}\label{eq:angular}
 W(T)=\sigma\left(\frac{\cos T-\cos r}{\cos(r/2)-\cos r}\right),
 \qquad \vartheta(T)=\arcsin(\sin T)\,W(T).
\end{equation}
The denominator is positive. This is a smooth $2\pi$-periodic
function: the support of $W$ avoids the nonsmooth points of
$\arcsin(\sin T)$. On the principal arc, it equals $T$ for
$|T|\le r/2$ and is zero for $|T|\ge r$. In addition,
\begin{equation}\label{eq:angularbounds}
 \vartheta'(T)\le1,\qquad
 \sin T\,\vartheta(T)\ge0,\qquad
 |\vartheta(T)|\le\frac\pi2|\sin T|.
\end{equation}
For the derivative bound, on the support write
$\vartheta'=W+T W'$. The second term is nonpositive, since $W$ decreases
away from zero on the principal arc, and $W\le1$. The global sign
and absolute-value estimates follow from the corresponding properties
of $\arcsin$ and $0\le W\le1$. Periodicity gives uniform bounds on
all derivatives.

We now define the torus Hamiltonian correction and primitive:
\begin{equation}\label{eq:primitive}
 H=\sum_{i\in M}h(A_i)\chi_i(A)\rho(S_i)\vartheta(T_i),
 \qquad \lambda_T=\beta-dH.
\end{equation}
Both are single-valued despite the angular notation. The identity
$d\lambda_T=\operatorname{Im}\Omega_T$ is immediate. The work is to
extend the one-form, since neither $\beta$ nor $H$ separately extends
across the coordinate divisors.

\section{Simultaneous cancellation at the boundary}

Fix $p\in Y$ and $I=I(p)$. Work in $U_I$ with $u_i=a_i$,
$v_i=b_i/Q_i$. Shrink to an open neighborhood of $p$ on which, for
every $i\in I$, the following conditions hold on its torus part:
\begin{enumerate}
\item $\chi_i=1$, while every correction centered at a mutable
      neighbor of $I$ is identically zero;
\item $A_i<-L-2\varepsilon$, so $h(A_i)=A_i+\kappa$;
\item the principal branch
      $\mathcal L_i=\operatorname{Log}(1-u_i v_i)$ is defined and
      small enough that $\rho(\operatorname{Re}\mathcal L_i)=1$
      and $\vartheta(\operatorname{Im}\mathcal L_i)
       =\operatorname{Im}\mathcal L_i$.
\end{enumerate}
These conditions can be imposed simultaneously. Every neighbor of a
zero coordinate is nonzero by independence, so the required strict
relative-norm inequalities hold near $p$. Also $u_i v_i=0$ at $p$,
and hence the logarithmic branch is defined near its value $1$.
The first condition makes neighboring corrections vanish on an open
neighborhood, including all their derivatives.

On the torus part,
$\mathcal L_i=S_i+\mathrm iT_i$, since $-q_i=1-u_i v_i$. The
terms of $\beta$ involving an index in $I$ are exactly
\[
 \sum_{i\in I}(A_i\,dT_i-S_i\,d\theta_i).
\]
There is no block with two indices in $I$. For each such index,
\begin{align}
 A_i\,dT_i-S_i\,d\theta_i-d((A_i+\kappa)T_i)
 &=-T_i\,dA_i-S_i\,d\theta_i-\kappa\,dT_i\notag\\
 &=-\operatorname{Im}\left(
        \mathcal L_i\frac{du_i}{u_i}+\kappa\,d\mathcal L_i
                        \right).\label{eq:primitivecancellation}
\end{align}
The potentially singular coefficient is removable. Namely,
\[
 \mathcal L_i/u_i=v_i F(u_i v_i),\qquad
 F(w)=\begin{cases}\operatorname{Log}(1-w)/w,&w\ne0,\\-1,&w=0.
             \end{cases}
\]
The power series $F(w)=-\sum_{n\ge0}w^n/(n+1)$ near zero shows that
$F$ is smooth there. Thus \eqref{eq:primitivecancellation} extends as
an actual smooth real covector, even if both $u_i$ and $v_i$ vanish.
At $u_i=0$ its value is
\begin{equation}\label{eq:boundaryprimitive}
 (1+\kappa)\operatorname{Im}(v_i\,du_i).
\end{equation}

All outside terms of $\beta$ involve units. Any correction not indexed
by $I$ or a neighbor of $I$ depends only on nonzero coordinates near
$p$, and is already smooth. Adding those regular terms to the
extensions \eqref{eq:primitivecancellation} defines a smooth local
one-form. It agrees with $\lambda_T$ on the dense torus. The local
one-forms therefore agree on overlaps and glue to a global $\lambda$.
Finally $d\lambda=\omega$ follows from equality on $T$ and continuity
of the two smooth two-forms. Together with Lemma~\ref{lem:nondegenerate},
this gives an exact symplectic structure on the actual manifold $Y$.

\section{The actual Liouville vector field and its boundary equations}

Nondegeneracy determines a unique smooth real vector field $Z$ by
$\iota_Z\omega=\lambda$. Its smoothness follows locally by smoothly
inverting the matrix of the nondegenerate two-form. We analyze this
field, rather than prescribing a separate model flow.

On the torus set $f_i(A)=h(A_i)\chi_i(A)\rho(S_i)$. Contracting
$\omega=\sum_{ij}B_{ij}\,dA_i\wedge d\theta_j$ gives
\begin{align}
 ZA_i&=A_i-h(A_i)\chi_i(A)\rho(S_i)\vartheta'(T_i)
                    &&(i\in M),\label{eq:radial}\\
 ZA_f&=A_f &&(f\in F),\label{eq:frozenradial}\\
 BZ\theta&=\partial_A H.\label{eq:angularvector}
\end{align}
There are no derivatives of $\chi_i$ in \eqref{eq:radial}; such
derivatives occur in the angular equation instead. The inverse in
\eqref{eq:angularvector} is the inverse of the full matrix $B$.

Since $0\le\chi_i,\rho\le1$ and $\vartheta'\le1$, the product
$\chi_i\rho\vartheta'$ is at most one, even where $\vartheta'$ is
negative. Equations \eqref{eq:hproperties}--\eqref{eq:hstrict} imply
\begin{equation}\label{eq:radialsign}
 \begin{array}{ll}
 A_i\ge-2\varepsilon &\Longrightarrow ZA_i=A_i,\\
 A_i<0 &\Longrightarrow ZA_i\le A_i-h(A_i)<0.
 \end{array}
\end{equation}
There is also a constant $c>0$, depending only on the fixed seed and
cutoffs, such that
\begin{equation}\label{eq:radialgrowth}
 |ZA_j|\le c|A_j|.
\end{equation}
For a mutable index this follows from
$|h(A_j)|\le|A_j|$ and a uniform bound for $|\vartheta'|$;
for a frozen index it follows from \eqref{eq:frozenradial}.

Each of these statements extends to the entire open locus $a_j\ne0$,
even when other coordinates vanish. On that locus both
$A_j=\log|a_j|$ and $ZA_j=\operatorname{Re}(Za_j/a_j)$ are smooth.
The equality or weak inequality holds on its dense intersection with
$T$ and hence everywhere. The strict negative conclusion follows
from the already strict scalar inequality $A_j-h(A_j)<0$.

To control zero coordinates, return to the localized neighborhood of
Section~5. There is a smooth real function $\nu_i$ on this neighborhood
for which
\begin{equation}\label{eq:nodedrift}
 Zu_i=(-\kappa+\mathrm i\nu_i)u_i,\qquad
 Z(u_i v_i)=(u_i v_i-1)\operatorname{Log}(1-u_i v_i).
\end{equation}
Here is why the angular coefficient is genuinely smooth. In
$B Z\theta=\partial_AH$, the singular corrections have the localized
form $(A_i+\kappa)T_i$. Their radial derivative is $T_i$, with no
remaining logarithmic divergence; neighboring corrections vanish
identically. The other derivatives are those of smooth functions of
the nonzero coordinates. Multiplication by the constant matrix
$B^{-1}$ gives a smooth real angular vector. On the dense torus its
$i$th component is the coefficient $\nu_i$ above. Moreover every
neighbor of $i\in I$ has uncorrected radial drift in this
neighborhood, so $ZS_i=S_i$, while
$ZT_i=\partial_{A_i}H=T_i$. Differentiating
$u_i v_i=1+q_i$ then proves the second equation. Smoothness extends
both equations across $u_i=0$.

Define near zero
\[
 E(w)=\begin{cases}
 (w-1)\operatorname{Log}(1-w)/w,&w\ne0,\\
 1,&w=0.
 \end{cases}
\]
Differentiating $u_i v_i$ in \eqref{eq:nodedrift} first where $u_i\ne0$
and then using continuity gives
\[
 Zv_i=(E(u_i v_i)+\kappa-\mathrm i\nu_i)v_i.
\]
Consequently the exact equation on the zero divisor is
\begin{equation}\label{eq:partnerzero}
 u_i=0\quad\Longrightarrow\quad
 Zv_i=(1+\kappa-\mathrm i\nu_i)v_i.
\end{equation}
No division by a zero original coordinate or a possibly zero partner
is used in the final identity.

\begin{lemma}\label{lem:zeroinvariance}
Along an integral curve on a connected time interval, the set of
vanishing original coordinates is constant. For an index in that set,
\begin{equation}\label{eq:partnernorm}
 |v_i(t)|^2=e^{2(1+\kappa)(t-s)}|v_i(s)|^2.
\end{equation}
\end{lemma}
\begin{proof}
Near any time at which $a_i=0$, \eqref{eq:nodedrift} gives the scalar
equation
\[
 \frac{d}{dt}|a_i|^2=-2\kappa|a_i|^2.
\]
The integrating factor $e^{2\kappa t}$ shows that a zero value
persists throughout a small connected time neighborhood, in both
directions. Thus the zero-time set is open in the curve's time
interval; it is also closed by continuity. Connectedness shows that
it is empty or the whole interval. Frozen coordinates are never zero
by definition of $Y$. This proves constancy of the whole zero set,
without assuming completeness beforehand.

For $i$ in that set, $Q_i$ is nonzero at every curve point by
Lemma~\ref{lem:zeros}, so $v_i=b_i/Q_i$ is smooth along the curve.
Taking the real part of \eqref{eq:partnerzero} gives
$(|v_i|^2)'=2(1+\kappa)|v_i|^2$, and the integrating-factor identity
gives \eqref{eq:partnernorm}. This includes $v_i(s)=0$.
\end{proof}

\section{Finite-time confinement and the complete flow}

\begin{lemma}\label{lem:confinement}
Fix $p\in Y$ and $T_0\ge0$. There is a compact subset $K(p,T_0)$ of
$Y$ such that any integral curve of $Z$ starting at $p$ remains in
$K(p,T_0)$ at every time in its interval of definition with
$|t|\le T_0$.
\end{lemma}
\begin{proof}
Let $I=I(p)$. By Lemma~\ref{lem:zeroinvariance}, precisely the
coordinates indexed by $I$ vanish along such a curve. For $j\notin I$,
apply Gronwall's inequality to \eqref{eq:radialgrowth}, in both time
directions, to obtain
\[
 |A_j(t)|\le |A_j(0)|e^{c|t|}.
\]
Let $R=\max_{j\notin I}|A_j(0)|e^{cT_0}$, with $R=0$ if the
index set is empty. Then every nonzero original coordinate satisfies
\begin{equation}\label{eq:modulusbounds}
 e^{-R}\le |a_j(t)|\le e^R\qquad(|t|\le T_0).
\end{equation}
The remaining originals are identically zero.

If a mutable $i$ is not in $I$, its partner is
$b_i=(P_i+Q_i)/a_i$. Each exchange monomial is a finite product of
nonnegative powers of the bounded originals, and the denominator is
bounded below by \eqref{eq:modulusbounds}; thus $b_i$ is uniformly
bounded on the specified finite time range. If $i\in I$, use
$b_i=Q_i v_i$. The monomial $Q_i$ involves only nonzero coordinates
and is uniformly bounded by \eqref{eq:modulusbounds}, while
\eqref{eq:partnernorm} gives
$|v_i(t)|\le e^{(1+\kappa)T_0}|v_i(0)|$. Hence these partners are
also uniformly bounded. This split between the two types of partner
is necessary; one cannot use $b_i=(P_i+Q_i)/a_i$ at $i\in I$.

Choose common finite bounds $A_0,B_0$ for all originals and partners,
and put $\eta=e^{-R}>0$. The coordinate box
\[
 \left\{q\in Y:|a_j(q)|\le A_0,\ |b_i(q)|\le B_0,
                  \ |a_f(q)|\ge\eta\ (f\in F)\right\}
\]
contains every such curve segment. Its image in
$\mathbb C^E\times\mathbb C^M$ is closed and bounded: the exchange
equations and the displayed inequalities are closed conditions, and
the positive frozen lower bound implies all the nonvanishing
conditions defining $Y$. Finite-dimensional compactness therefore
proves the claim in the inherited topology.
\end{proof}

\begin{proposition}\label{prop:complete}
The vector field $Z$ has a unique smooth complete flow
$\varphi:\mathbb R\times Y\to Y$.
\end{proposition}
\begin{proof}
A smooth vector field has local solutions, uniqueness, and smooth
dependence in the node charts. On a compact set, finitely many local
existence neighborhoods yield a common positive lifetime for solutions
starting at any point of that set.

Suppose a maximal solution through $p$ had a finite right endpoint
$b$. Lemma~\ref{lem:confinement}, applied to a time bound larger than
$|b|$, puts its segment approaching $b$ in a fixed compact set. Let
$\tau>0$ be a common two-sided lifetime for initial points of that
set. Start a local solution at a time $s<b$ with $b-s<\tau/2$.
Uniqueness identifies it with the old solution on their overlap,
and it extends the old solution beyond $b$, a contradiction. The
same argument at a finite left endpoint proves existence for all
real times.

Set $\varphi_t(p)$ equal to this all-real solution at time $t$.
Uniqueness yields $\varphi_0=\mathrm{id}$ and
$\varphi_{s+t}=\varphi_s\circ\varphi_t$. Local smooth dependence,
continued along a finite orbit segment by finitely many local flow
maps, gives joint smoothness in $(t,p)$. The resulting flow is the
flow of the already constructed field $Z$, and its orbit derivative
is $Z(\varphi_t(p))$.
\end{proof}

\section{Compactness and positive invariance of the coordinate cap}

\begin{lemma}\label{lem:capcompact}
The set $C$ is compact. At a point of $C$, the set
$J=\{i\in M:|a_i|<1\}$ is independent, and every neighbor of an
index in $J$ has modulus one.
\end{lemma}
\begin{proof}
All original coordinates have modulus at most one. If $|a_i|<1$,
the cap equation gives $b_i=0$. If $|a_i|=1$, the exchange equation
gives $|b_i|\le|P_i|+|Q_i|\le2$. The cap is thus bounded in the
ambient coordinate space. Its image there is closed, being defined
by the exchange equations, the frozen unit-modulus equations, and
the closed conditions in \eqref{eq:cap}. The frozen equations already
imply nonvanishing. Hence it is compact.

For independence, if the mutable graph induced on $J$ had an edge,
choose a source in a nontrivial connected component. At that source
$i$ the partner vanishes, so $P_i+Q_i=0$. Every incoming factor has
modulus one, while at least one outgoing factor has modulus less
than one. Thus $|Q_i|=1$ and $|P_i|<1$, a contradiction. Therefore
$J$ is independent. All mutable coordinates outside $J$ have modulus
one by definition, and all frozen coordinates have modulus one by
the cap equations. This proves the neighbor assertion.
\end{proof}

We record a character equation that will be used in both inclusions.
Suppose $a_i\ne0$ and every full-row neighbor of $i$ has modulus
one. Then $Q_i\ne0$, $|q_i|=1$, and $S_i=0$. Put
\[
 g_i(x)=h(x)\prod_{\substack{j\in M\\B_{ij}\ne0}}\chi(x),
 \qquad c_i(x)=g_i'(x).
\]
In this formula the neighbors have logarithmic radius zero. The
factor $\rho(S_i)=\rho(0)=1$ has no derivative with respect to $A_i$,
since $B_{ii}=0$. The full angular equation gives the exact identity
\begin{equation}\label{eq:characterode}
 \frac{d}{dt}q_i
   =\mathrm i\,q_i\,c_i(A_i)\,
                    \vartheta(\arg(-q_i)).
\end{equation}
Indeed terms centered at a mutable neighbor of $i$ have both $h(0)$
and $h'(0)$ equal to zero; terms at nonneighbors have no $A_i$
dependence. Only the term centered at $i$ remains in
$\partial_{A_i}H$. The same calculation is valid on simultaneous
zero strata elsewhere: only row-support coordinates enter this
character and they are nonzero.

There is a useful branch-free uniqueness estimate. From
\eqref{eq:angularbounds} and the polar sine identity,
\begin{equation}\label{eq:characterestimate}
 |z|\,|\vartheta(\arg(-z))|
 \le\frac\pi2|\operatorname{Im}z|
 \le\frac\pi2|z+1|.
\end{equation}
Thus a solution of \eqref{eq:characterode} that equals $-1$ at one
time remains $-1$ on any finite time interval on which the
coefficient is continuous. To see this directly, bound $|c_i(A_i)|$
on that interval and apply Gronwall to $|q_i+1|$ using
\eqref{eq:characterestimate}. Reversing the time parameter gives
the same uniqueness in the backward direction. This argument does
not differentiate a discontinuous choice of argument.

\begin{proposition}\label{prop:capinvariant}
The cap is positively invariant under the complete flow $\varphi$.
Consequently $C\subseteq\operatorname{Skel}(\varphi)$.
\end{proposition}
\begin{proof}
Let $p\in C$. First, any original coordinate initially of modulus
one stays of modulus one: it is nonzero throughout the orbit by
Lemma~\ref{lem:zeroinvariance}, and its logarithmic radius is
identically zero by \eqref{eq:radialgrowth} and Gronwall. Every
initially nonzero coordinate of modulus less than one stays strictly
inside the disk, since \eqref{eq:radialsign} makes its negative
logarithmic radius decrease. These statements use the vector field
equations and the initial values; they do not assume cap membership
at later times.

Consider an initially nonzero strict-interior mutable coordinate
$a_i$. Lemma~\ref{lem:capcompact} gives unit modulus at all its
neighbors, and the preceding zero-log argument propagates that
condition along the orbit. The cap equation initially gives
$b_i=0$, so the normalized exchange equation gives $q_i=-1$.
Equation~\eqref{eq:characterode} and the uniqueness estimate imply
$q_i(t)=-1$ for every $t\ge0$. Since $a_i(t)$ and $Q_i(t)$ remain
nonzero, the exchange equation then gives $b_i(t)=0$.

If $a_i(p)=0$, it stays zero by
Lemma~\ref{lem:zeroinvariance}. Its initial partner is zero by the
cap condition, so \eqref{eq:partnernorm} gives $v_i(t)=0$ and hence
$b_i(t)=0$. At indices with initial modulus one, no partner
restriction is needed because $1-|a_i(t)|^2=0$. Frozen coordinates
retain modulus one as well. All the conditions of \eqref{eq:cap}
therefore hold for all positive times.

The entire forward orbit of $p$ lies in the compact set $C$.
Its closure is a closed subset of $C$, so is compact. This proves
the stated inclusion using the independent definition
\eqref{eq:trapped}.
\end{proof}

\section{Every trapped orbit lies in the cap}

Let $p\in\operatorname{Skel}(\varphi)$ and write
$\gamma(t)=\varphi_t(p)$. Denote its compact forward-orbit closure
by $K$. In particular, all original and partner coordinates are
bounded on $K$, and every frozen modulus is bounded below by a
positive constant on $K$.

\subsection{Radial restrictions and independent limiting directions}

For a frozen coordinate the exact solution of
\eqref{eq:frozenradial} is $A_f(t)=e^tA_f(0)$. A positive initial
value makes $|a_f|$ unbounded, while a negative value makes it tend
to zero, contradicting the positive lower bound on $K$. Thus
$|a_f(p)|=1$.

For any mutable coordinate with $A_i(0)>0$, equation
\eqref{eq:radialsign} gives the same positive exponential law and
contradicts compactness. Hence every original modulus is at most
one initially. For $0<|a_i(p)|<1$, its logarithmic radius remains
negative and strictly decreases. It tends to $-\infty$: if it had a
finite negative limit $\ell$, the continuous strict bound
$A-h(A)<0$ would give $A_i'\le-c_0<0$ once $A_i$ is sufficiently
close to $\ell$, contradicting convergence to $\ell$.
Consequently
\begin{equation}\label{eq:interiorlimit}
 |a_i(p)|<1\quad\Longrightarrow\quad |a_i(\gamma(t))|\longrightarrow0,
\end{equation}
including an initially zero coordinate, which stays zero.

Set $J=\{j\in E:|a_j(p)|<1\}$. No frozen index lies in $J$.
Choose a sequence of times tending to $+\infty$ and a convergent
subsequence of their points in $K$. By \eqref{eq:interiorlimit},
its limit has every $J$-coordinate zero. The zero set at the limit
is independent by Lemma~\ref{lem:zeros}; therefore $J$ is
independent. Every neighbor of an index in $J$ has initial modulus
one and hence modulus one for the entire orbit, by the zero-log law.
This compact-limit argument includes arbitrary simultaneous-zero
strata; it does not restrict the orbit closure to a single open
node chart.

\subsection{Nonzero interior coordinates: angular rigidity}

Fix a mutable $i$ with $0<|a_i(p)|<1$. All its neighbors have
modulus one, so $|Q_i(\gamma(t))|=1$ and the genuine character
$q_i=P_i/Q_i$ satisfies \eqref{eq:characterode} throughout the
orbit. Boundedness of $b_i$ on $K$ and $a_i(t)\to0$ give
\begin{equation}\label{eq:qlimit}
 q_i(t)+1=\frac{a_i(t)b_i(t)}{Q_i(t)}\longrightarrow0.
\end{equation}
Eventually $A_i(t)<-L-2\varepsilon$. In this deep region the
neighbor radii are zero, so each dominance factor is one and
$g_i(A_i)=A_i+\kappa$. Thus $c_i(A_i)=1$ at all sufficiently large
times.

To conclude rigidity without choosing angle lifts, differentiate
the squared distance to the equilibrium:
\begin{equation}\label{eq:distancemonotone}
 \frac{d}{dt}|q_i+1|^2
 =2c_i(A_i)\bigl(-\operatorname{Im}q_i\bigr)
                      \vartheta(\arg(-q_i)).
\end{equation}
The last two factors have nonnegative product, because
$-\operatorname{Im}q_i=|q_i|\sin(\arg(-q_i))$ and
$\sin T\,\vartheta(T)\ge0$. Therefore $|q_i+1|^2$ is
nondecreasing for all sufficiently large times. By
\eqref{eq:qlimit} it tends to zero, so it must already be identically
zero on that tail.

Apply the finite-interval backward uniqueness following
\eqref{eq:characterestimate} to propagate $q_i=-1$ from the tail
back to the initial time. This step is necessary: in the transition
region the derivative $c_i(A_i)$ need not be nonnegative, and the
dominance factor need not be one. We use only continuity of that
coefficient on each finite interval. Finally $a_i(p)\ne0$ and
$Q_i(p)\ne0$ in the exchange equation give $b_i(p)=0$.

\subsection{Initially zero coordinates: normalized-partner escape}

Now let $a_i(p)=0$. Its entire forward orbit lies in the closed
zero divisor $D_i=\{a_i=0\}$. On this whole divisor $Q_i$ is
nonzero: at each point, Lemma~\ref{lem:zeros} says that every
neighbor of $i$ is nonzero. Thus $b_i/Q_i$ is continuous at every
point of $D_i$, including points with additional zero coordinates
away from the neighbors of $i$.

The intersection $K\cap D_i$ is compact, so $|b_i/Q_i|$ is bounded
there. Formula~\eqref{eq:partnernorm}, however, gives
\[
 |b_i(\gamma(t))/Q_i(\gamma(t))|^2
 =e^{2(1+\kappa)t}|b_i(p)/Q_i(p)|^2.
\]
Since $1+\kappa>0$, boundedness forces the initial normalized
partner to be zero, and hence $b_i(p)=0$. It would not suffice here
to assume that a normalization is bounded in one chosen open
chart; compactness is applied on the entire closed zero slice.

\begin{proposition}\label{prop:reverse}
$\operatorname{Skel}(\varphi)\subseteq C$.
\end{proposition}
\begin{proof}
The radial argument gives $|a_f(p)|=1$ for frozen indices and
$|a_i(p)|\le1$ for mutable indices. If $|a_i(p)|<1$, the two
preceding cases give $b_i(p)=0$, according as $a_i(p)$ is nonzero
or zero. If $|a_i(p)|=1$, the product
$(1-|a_i(p)|^2)b_i(p)$ vanishes automatically. These are exactly the
conditions defining $C$.
\end{proof}

Propositions~\ref{prop:capinvariant} and~\ref{prop:reverse} establish
the literal set equality for the same complete flow. Together with
the manifold construction and $d\lambda=\omega$, this proves
Theorem~\ref{thm:main}.

\section{Correspondence with the formal statement}

The point type in the formalization is \texttt{Seed.Point}: the
coordinate subtype \eqref{eq:exchange}, with its inherited topology.
\texttt{SmoothPresentation} requires the actual coordinate inclusion
to be a smooth immersion. The real tangent spaces, manifold
derivatives, exterior derivative, and flow are native Mathlib objects
used through the pinned \texttt{FsFltz.Geometry} interfaces. The
construction is therefore on this coordinate space, not on an
unrelated carrier transported by an arbitrary point equivalence.

The final theorem is
\begin{center}
\texttt{ClusterSkeleton.FullSkeleton.firstGeometricTarget}.
\end{center}
Its quantified seed hypotheses are precisely skew-symmetry, mutable
acyclicity, and determinant one for the full exchange matrix. The
parameterized construction accepts \eqref{eq:parameters}, and the
unparameterized existence result chooses $(1,3,1)$. None of these
three parameters fixes a frozen coordinate.

The two sides of the equality are defined independently:
\texttt{coordinateCap} is \eqref{eq:cap}, while
\texttt{trappedSet} is compactness of the closure of the forward
orbit of a Mathlib \texttt{Flow}. The final witnesses contain the
global primitive, its Liouville vector field, and the constructed
complete flow. Their contraction equation, all-real existence,
joint smoothness, and the skeleton equality are proved properties;
they are not supplied assumptions in the final existence theorem.

\begin{longtable}{>{\raggedright\arraybackslash}p{0.33\linewidth}>{\raggedright\arraybackslash}p{0.61\linewidth}}
\toprule
Mathematical step & Main source files or final declaration\\
\midrule
\endhead
Exchange equations and zero-set independence &
\texttt{Seed.lean}, \texttt{ZeroSet.lean}\\[3pt]
Integral node coordinates and native smooth atlas &
\texttt{IndependentLaurentCoordinates.lean},
\texttt{IndependentNodeChart.lean}, \texttt{SmoothAtlas.lean}\\[3pt]
Density and full local two-form &
\texttt{InitialTorusDensity.lean}, \texttt{NodeTwoForm.lean}\\[3pt]
Residual-block nondegeneracy and global symplectic form &
\texttt{NodeNondegeneracy.lean}, \texttt{GlobalTwoForm.lean},
\texttt{GlobalSymplectic.lean}\\[3pt]
Explicit cutoffs and simultaneous primitive cancellation &
\texttt{PrimitiveCutoffs.lean}, \texttt{DominanceCutoff.lean},
\texttt{NodePrimitiveCancellation.lean}, \texttt{LocalPrimitive.lean}\\[3pt]
Global primitive and its actual field &
\texttt{GlobalPrimitive.lean}, \texttt{GlobalLiouvilleField.lean}\\[3pt]
Zero-stratum invariance and complex partner dynamics &
\texttt{ZeroStratumInvariance.lean}, \texttt{LocalAngularDynamics.lean},
\texttt{ZeroStratumPartnerDynamics.lean}\\[3pt]
Finite-time compact confinement and all-real native flow &
\texttt{FiniteTimeConfinement.lean}, \texttt{NativeCompleteFlow.lean}\\[3pt]
Cap compactness, interior partner propagation, and invariance &
\texttt{CapSet.lean}, \texttt{InteriorPartnerPropagation.lean},
\texttt{CapForwardInvariance.lean}\\[3pt]
Compact-orbit rigidity and zero-partner escape &
\texttt{CompactInteriorRigidity.lean}, \texttt{ZeroStratumEscape.lean}\\[3pt]
The two inclusions for the constructed flow &
\texttt{ConstructedFlowCapInclusion.lean},
\texttt{ConstructedFlowReverseInclusion.lean}\\[3pt]
The literal equality and full existence theorem &
\texttt{FullSkeleton.trappedSet\_eq\_coordinateCap},
\texttt{FullSkeleton.liouvilleSkeletonTarget},
\texttt{FullSkeleton.firstGeometricTarget}\\
\bottomrule
\end{longtable}

The accompanying verification records concern this geometric theorem.
The proof above uses no conclusion from homological mirror symmetry or
from a finite-type or cocore construction. Its scope is exactly the
real Liouville structure and full trapped-set equality stated in
Theorem~\ref{thm:main}.

\paragraph{Sources.}
The motivating geometric argument is the Liouville and skeleton portion
of the supplied \emph{acyclic\_cluster\_HMS} notes, Section~2.
This exposition has been rewritten to match the formalized real-geometric
construction; it does not reproduce the unrelated parts of those notes.
The Lean project uses Mathlib and Bohan Fang's
\href{https://github.com/bohanfang/fs-fltz-minimal}{\texttt{fs-fltz-minimal}}
as pinned dependencies. Their exact revisions and the original manuscript
hash are recorded in \texttt{docs/sources.json}.

\end{document}
